\section{Trabalhos relacionados: revisões existentes}
\label{sec:relacionados}

A \emph{string} secundária devolveu \nSecundariaTotal{} registos
(\nSecundariaUnicos{} únicos). Destes, e das revisões excluídas por CE4 na triagem
da \emph{string} principal, selecionaram-se as quatro mais diretamente
relacionadas com a segurança do MQTT; a quinta, de Hintaw et al., foi acrescentada
por \emph{snowballing} (Secção~\ref{sec:selecao}).

\begin{itemize}[leftmargin=*]
  \item \textbf{Perrone et al.\ (2017)}~\cite{perrone2017mirai}, escrito após os
    ataques da botnet Mirai, descreve as soluções de segurança habitualmente usadas
    em instalações MQTT reais e fecha com um levantamento de temas de investigação
    que os próprios autores descrevem como conciso e não exaustivo.
  \item \textbf{Nebbione e Calzarossa (2020)}~\cite{nebbione2020security} revêem a
    segurança dos protocolos da camada de aplicação da IoT, incluindo o MQTT, as
    vulnerabilidades publicadas (CVE) e boas práticas de mitigação.
  \item \textbf{Hintaw et al.\ (2023)}~\cite{hintaw2023mqtt} organizam os problemas
    de segurança do MQTT em negação de serviço, falsificação de identidade,
    divulgação de informação, elevação de privilégios e adulteração, e apontam o
    custo da segurança em dispositivos a bateria como desafio em aberto.
  \item \textbf{Lakshminarayana et al.\ (2024)}~\cite{lakshminarayana2024securing}
    classificam as vulnerabilidades e os ataques ao MQTT e discutem as defesas, com
    ênfase na identificação de vulnerabilidades e na deteção de ataques.
  \item \textbf{Laghari et al.\ (2024)}~\cite{laghari2024securing} apresentam uma
    taxonomia de ataques ao ecossistema MQTT e apontam às revisões anteriores uma
    discussão ampla mas superficial.
\end{itemize}

\begin{table}[ht]
  \centering
  \caption{Comparação das revisões relacionadas com a presente revisão.
    \Sim{} abordado; \Parc{} parcialmente; \Nr{} não reportado. $^{*}$Classificada
    a partir do resumo: o texto integral não estava acessível.}
  \label{tab:revisoes}
  \small
  \setlength{\tabcolsep}{3pt}
  \begin{tabular}{@{}lcC{1.45cm}C{1.75cm}*{4}{C{1.45cm}}@{}}
    \toprule
    \textbf{Revisão} & \textbf{Ano}
      & {\scriptsize\textbf{Só MQTT}}
      & {\scriptsize\textbf{Taxonomia de ataques}}
      & {\scriptsize\textbf{Revê defesas}}
      & {\scriptsize\textbf{Compara por propriedade}}
      & {\scriptsize\textbf{Protocolo RSL explícito}}
      & {\scriptsize\textbf{Propõe arquitetura}} \\
    \midrule
    Perrone et al.~\cite{perrone2017mirai}              & 2017 & \Sim  & \Nr   & \Sim & \Nr & \Nr & \Nr \\
    Nebbione e Calzarossa~\cite{nebbione2020security}   & 2020 & \Parc & \Sim  & \Sim & \Nr & \Nr & \Nr \\
    Hintaw et al.~\cite{hintaw2023mqtt}$^{*}$          & 2023 & \Sim  & \Parc & \Sim & \Nr & \Nr & \Nr \\
    Lakshminarayana et al.~\cite{lakshminarayana2024securing} & 2024 & \Sim & \Sim & \Sim & \Nr & \Nr & \Nr \\
    Laghari et al.~\cite{laghari2024securing}           & 2024 & \Sim  & \Sim  & \Sim & \Nr & \Parc & \Nr \\
    \midrule
    \rowcolor{fundotab}
    \textbf{Este trabalho}                              & 2026 & \Sim  & \Parc & \Sim & \Sim & \Sim & \Sim \\
    \bottomrule
  \end{tabular}
\end{table}

As revisões existentes organizam-se em torno dos \emph{ataques}
(Tabela~\ref{tab:revisoes}). Só Laghari et al.\ descrevem a pesquisa e a seleção,
e apenas das revisões que comparam, com uma única palavra-chave em três
fontes~\cite[§I]{laghari2024securing}. Esta revisão distingue-se por relatar um
protocolo explícito e reprodutível, por comparar os estudos segundo as propriedades que garantem à
\emph{mensagem} (incluindo validade temporal, anti-repetição e não repúdio) e a
confiança que depositam no \emph{broker}, por cobrir trabalhos até 2026 e por
ligar as lacunas encontradas a uma proposta concreta.
